
\section{Hermite 計量・Kähler 計量と正の直線束（復習）}\label{sec:herm}

\subsection{Hermite 計量と Kähler 形式}
複素多様体 $X$ の\emph{Hermite 計量}とは、$T^{1,0}X$ 上の Hermite 内積 $h=\sum h_{j\bar k}\,dz_j\otimes d\bar z_k$ を滑らかに与えることである。同値に、$J$-不変な Riemann 計量 $g=\mathrm{Re}\,h$ と、実 $(1,1)$ 形式
\[
\omega=g(J\cdot,\cdot)=i\sum_{j,k}g_{j\bar k}\,dz_j\wedge d\bar z_k
\]
（\emph{基本形式}）を与えることでもある（第2巻 \S6--7）。$\C^n$ の標準計量では $g_{j\bar k}=\frac12\delta_{jk}$、$\omega=\frac i2\sum dz_j\wedge d\bar z_j=\sum dx_j\wedge dy_j$ である。$\omega^n/n!$ は体積形式で、$X$ には $\omega^n/n!$ で向きが入る。

\begin{defn}[Kähler]
$d\omega=0$ のとき $g$ を\emph{Kähler 計量}という。同値な条件（第2巻、定理 Kähler 条件の同値）：各点のまわりで $\omega=i\partial\bar\partial\varphi$（$\varphi$ は実関数、Kähler ポテンシャル）と書ける。また、各点 $p$ のまわりに正則座標で $g_{j\bar k}=\frac12\delta_{jk}+O(|z|^2)$ となるもの（\emph{Kähler 正規座標}）がとれる\tC\ （Huybrechts, \emph{Complex Geometry}, 第1章・第3章；細部は文献）。
\end{defn}
最後の性質は本巻で Kähler 恒等式を証明するときに使う（\S\ref{sec:kid}）。複素 1 次元では $d\omega$ は 3 形式なので、すべての Hermite 計量が Kähler である。

\begin{ex}\label{ex:kahlerlist}
(1) $\C^n$、複素トーラス $\C^n/\Gamma$（$\omega=\frac i2\sum dz_j\wedge d\bar z_j$）。(2) $\CP^n$ の Fubini--Study 計量 $\omega_{FS}=\frac i{2\pi}\partial\bar\partial\log|Z|^2$（$\int_{\CP^n}\omega_{FS}^n=1$）。(3) Kähler 多様体の複素部分多様体（誘導計量）。特に滑らかな射影多様体はすべて Kähler。(4) 反例：Hopf 曲面 $(\C^2\setminus0)/\langle z\mapsto2z\rangle$（微分同相型 $S^1\times S^3$）はコンパクトだが Kähler でない（第2巻：$b_1=1$ が奇数、または $H^2_{dR}=0$ で $[\omega]=0$）。
\end{ex}

\subsection{正の直線束}
$(L,h)$ を Hermite 直線束とする。局所正則枠 $e$ について $h:=h(e,e)>0$ とおくと、Chern 接続の曲率は $\Theta_h=\bar\partial\partial\log h=-\partial\bar\partial\log h$、第 1 Chern 類は $c_1(L)=\big[\frac i{2\pi}\Theta_h\big]$ である（第2巻 \S6）。たとえば $h=e^{-\varphi}$ なら $i\Theta_h=i\partial\bar\partial\varphi$ である。

\begin{defn}[正の直線束]\label{def:positive}
正則直線束 $L$ が\emph{正}（positive）であるとは、Hermite 計量 $h$ で、曲率形式
\[
i\Theta_h=i\sum_{j,k}R_{j\bar k}\,dz_j\wedge d\bar z_k,\qquad (R_{j\bar k})\ \text{が各点で正定値 Hermite 行列}
\]
となるものが存在することをいう。$L$ が\emph{負}とは $L^{-1}$ が正のこと。局所的には「$-\log h$ が強多重劣調和」と言い換えられる。
\end{defn}

\begin{prop}\label{prop:posKahler}\tP
コンパクト複素多様体 $X$ に正の直線束 $(L,h)$ があれば、$\omega:=i\Theta_h$ は Kähler 形式で、$[\omega]=2\pi c_1(L)$。特に $X$ は Kähler である。
\end{prop}
\begin{proof}
局所的に $\omega=i\partial\bar\partial\varphi$（$\varphi=-\log h$）で、$\varphi$ は実関数。$d(i\partial\bar\partial\varphi)=i(\partial+\bar\partial)\partial\bar\partial\varphi=0$。正定値性は仮定。$\frac i{2\pi}\Theta_h$ の類が $c_1(L)$ であることは定義。
\end{proof}
したがって「$X$ は Kähler でない」という事実は「$X$ に正の直線束はない」を意味する。Hopf 曲面はその例（問 \ref{p:hopfnopos}）。

\begin{ex}[$\Oc(1)$ は正]\label{ex:O1pos}
$\CP^n$ の座標 $[Z_0:\dots:Z_n]$ に対し、タウトロジー直線束 $\Oc(-1)\subset\CP^n\times\C^{n+1}$ に $\C^{n+1}$ の標準計量を制限する。枠 $e_0=(1,z_1,\dots,z_n)$（$U_0$ 上、$z_j=Z_j/Z_0$）について $h(e_0,e_0)=1+|z|^2$。第2巻の例のとおり、双対 $\Oc(1)$ の計量では $h_{\Oc(1)}=(1+|z|^2)^{-1}$ で、
\[
\frac i{2\pi}\Theta_{\Oc(1)}=\frac i{2\pi}\partial\bar\partial\log(1+|z|^2)=\omega_{FS}>0,\qquad c_1(\Oc(1))=[\omega_{FS}].
\]
したがって $\Oc(1)$ は正、$\Oc(m)$ は $m>0$ で正、$m<0$ で負、$m=0$ で正でも負でもない。滑らかな部分多様体 $Y\subset\CP^N$ では $\Oc_Y(1)=\Oc(1)|_Y$ が正（問 \ref{p:fsposit}）。
\end{ex}

\begin{ex}[曲線上の正値性 $\iff$ 次数 $>0$]
Riemann 面 $X$ では $\int_X\frac i{2\pi}\Theta_h=\deg L$（第4巻）。$L$ が正なら曲率形式は正の面積形式なので $\deg L>0$。逆も成り立つ（問 \ref{p:posdeg}：曲率の密度が一定になるよう計量を取り替える）。したがって 1 次元では「正 $\iff\deg>0$」。第4巻の定理（$\deg L>0$ なら $H^1(K\otimes L)=0$ など）は、本巻の高次元の定理の $n=1$ の場合に当たる。
\end{ex}

\begin{keybox}{要点（\S\ref{sec:herm}）}
$\omega=g(J\cdot,\cdot)=i\sum g_{j\bar k}dz_j\wedge d\bar z_k$。Kähler $\iff d\omega=0\iff$ 局所的に $\omega=i\partial\bar\partial\varphi$ $\iff$ Kähler 正規座標がある。直線束 $L$ が正とは、計量 $h$ で $i\Theta_h>0$ となるもの（$h=e^{-\varphi}$ なら $\varphi$ が強多重劣調和）。コンパクト $X$ に正の $L$ があれば $\omega=i\Theta_h$ が Kähler 形式（$X$ の計量を自由に選べるのが Kodaira の方法の要）。$\Oc(1)$ は $\CP^n$ で正。
\end{keybox}

\begin{prob}{基本}\label{p:fsposit}
(a) $\CP^n$ の $U_0$ で $\omega_{FS}=\frac i{2\pi}\partial\bar\partial\log(1+|z|^2)$ の $z=0$ における値が $\frac i{2\pi}\sum dz_j\wedge d\bar z_j$ であることを確かめ、正定値であることを示せ。(b) $U(n+1)$ の作用による不変性から、$\omega_{FS}$ は $\CP^n$ の至る所で正定値であることを説明せよ。(c) $Y\subset\CP^N$ を複素部分多様体とするとき、$\Oc(1)|_Y$ が正であることを示せ。
\end{prob}
\begin{prob}{標準}\label{p:posdeg}
$X$ をコンパクト Riemann 面、$L$ を $\deg L>0$ の正則直線束とする。任意の Hermite 計量 $h_0$ から出発し、$h=e^{-\psi}h_0$ とおく。曲率形式の密度が一定（$=\deg L\cdot\mathrm{vol}^{-1}$）になる $\psi$ が存在することを、Poisson 方程式 $\Delta\psi=f$（$f$ は与えられた関数で $\int f\,dA=0$ なら解ける。第7巻、または Hodge の定理 \S\ref{sec:hodge}）から導き、$L$ が正であることを示せ。
\end{prob}
\begin{prob}{標準}\label{p:hopfnopos}
Hopf 曲面 $X$ には正の直線束がなく、したがって射影的でない（どの $\CP^N$ にも複素部分多様体として埋め込めない）ことを示せ。
\end{prob}

